\documentclass{article}
\usepackage{graphicx}
\usepackage{amsmath}
\usepackage{amssymb}
\usepackage{amsfonts}
\usepackage{hyperref}
\usepackage{exercises}
\usepackage{indentfirst}
\usepackage{amsthm}
\usepackage{csquotes}
\usepackage{yhmath}
\usepackage{comment}
\usepackage{placeins}
\usepackage{tikz}
\newcommand*{\oldneg}{\mathord{\sim}}
\newcommand{\Z}{\mathbb{Z}}
\newcommand{\N}{\mathbb{N}}
\let\oldemptyset\emptyset

\theoremstyle{definition}
\newtheorem{definition}{Definition}

\setlength{\parindent}{.5in}
\setlength{\parskip}{12pt}
\title{Chapter 5 Answers}
\begin{document}
\date{}
\maketitle

\section*{Problem 1}
(a): Consider $\omega_n$, that is, the set of all $x$ such that $r\left(n,x\right)\in T_0$. In the last chapter, we showed that $T_0$ and the relation 
$g\left(x\right)=y$ are both recursively enumerable. Hence $r\left(n,x\right)\in T_0$ iff $g\left(x\right)=x_0\wedge nx_0\in T_0$. Now, the relation 
$r\left(x,y\right)=z$ is recursively enumerable, as it is expressed by: \[ \exists x_0\left(g\left(x\right)=x_0\wedge yx_0=z\right) \]
and thus the relation $r\left(n,x\right)\in T_0$ is expressed by the formula: \[ \exists z\left(r\left(n,x\right)=z\wedge z\in T_0\right) \]


(b): Suppose that $A$ is a recursively enumerable set. Hence there is some predicate $P$ of $\left(U\right)$ such that $Hx\in T$ iff $x\in A$. If $p$ is the Gödel number of 
$P$, $Px\in T$ iff $px_0\in T_0$ iff $r\left(p,x\right)\in T_0$ iff $x\in\omega_p$, thus proving that $A=\omega_p$.

\section*{Exercises 1-3}
I could not solve any of these.

\section*{Problem 2}
Consider a recursively enumerable relation $R\left(x,y\right)$. Hence, in $\left(U\right)$, it is represented by some predicate $H$, with a Gödel number of $h$. 
Hence $R(x,y)$ holds iff $hxcy\in T_0$ for $c$ the Gödel number of the comma. Take $t(y)$ to be $hy_0c$, which is recursive as $g(x)=x_0$ is recursively enumerable. 
Hence $x\in\omega_{t(i)}$ iff $r\left(hi_0c,x\right)\in T_0$ iff $hi_0cx\in T_0$ iff $R\left(i,x\right)$.


$Note$: The above is essentially the answer that Smullyan gave, but it is incorrect. It proves that $x\in\omega_{t(i)}$ iff $R(i,x)$, but that was not the theorem! It was 
to prove that $x\in\omega_{t(i)}$ iff $R(x,i)$. Below is a correct proof:


We suppose that $R(x,y)$ is recursively enumerable. Let $R_1(x,y)$ be the relation $R(y,x)$. Hence $R_1(x,y)$ is recursively enumerable. Thus 
$R(x,y)$ is represented in $(U)$ by some predicate $H$ with Gödel number $h$, and $R_1(x,y)$ is represented in $(U)$ by some predicate $P$ with Gödel number $p$. Hence 
\[ px_0cy_0\in T_0\textrm{ iff } hy_0cx_0\in T_0 \] for $x_0,y_0,c$ the Gödel numbers for $x,y,$ and the comma, respectively. Let $t(y)$ be the relation 
$py_0c$. Thus $x\in\omega_{t(i)}$ iff $r(pi_0c,x)\in T_0$ iff $pi_0cx_0\in T_0$ iff $hx_0ci_0\in T_0$ iff $R(x,i)$, meaning 
$x\in\omega_{t(i)}$ iff $R(x,i)$, thus completing the proof.

\section*{Problem 3}
Consider some recursively enumerable relation $R\left(x_1,...,x_m,y_1,...,y_n\right)$. Hence there is a recursively enumerable relation $M(x,y)$ such that 
$R(x_1,...,x_m,y_1,...,y_n)$ iff $M\left(\delta_m\left(x_1,...,x_m\right),\delta_n\left(y_1,...,y_n\right)\right)$. By the Simple Iteration Theorem, there is a function 
$t(i)$ such that $x\in\omega_{t(i)}$ iff $M(x,i)$. Hence, take 
$\phi(i_1,...,i_n)$ to be $t(\delta_n(i_1,...,i_n)$. Hence $M(\delta_m(x_1,...,x_m),\delta_n(i_1,...,i_n)$ iff $\delta_m(x_1,...,x_m)\in \omega_{\phi(i_1,...,i_n)}$ iff
$R_{\phi(i_1,...,i_n)}(x_1,...,x_m)$. This is true as we recall that $R^n_i(x_1,...,x_n)$ was defined to be $\delta_n(x_1,...,x_n)\in\omega_i$. 

\section*{Problem 4}
We consider the enumeration $\omega_1,\omega_2,...,\omega_n,...$ and any other recursively enumerable indexing $A_1,A_2,...,A_n,...$. Hence 
$x\in A_y$ is recursively enumerable, which means there is a recursive function $t(i)$ such that $x\in\omega_{t(i)}$ iff $x\in A_i$, thus completing the proof.

\section*{Problem 5}
Define $S(x_1,...,x_m,y,y_1,...,y_n)$ iff $R_y(x_1,...,x_m,y_1,...,y_n)$, which is recursively enumerable because $R_y(x_1,...,x_m,y_1,...,y_n)$ is recursively enumerable 
by the enumeration theorem. Hence, by the iteration theorem, there exists some $\phi(i,i_1,...,i_n)$ such that $R_{\phi(i,i_1,...,i_n)}(x_1,...,x_m)$ iff
$S(x_1,...,x_m,y,y_1,...,y_n)$ iff $R_i(x_1,...,x_m,i_1,...,i_n)$.

\section*{Problem 6}
Let $n$ be the number $t(m)$. Hence $R(x,n)$ iff $R_m(x,m)$ iff $x\in\omega_{n}$.


$Note$: To be clear, if $R(x,y)$ is recursively enumerable, then so is $R(x,t(y))$, as it is equivalent to the recursively enumerable relation:
\[ \exists z(t(y)=z\wedge R(x,z)) \] By the enumeration theorem, there then exists a relation $R_m(x,y)$ such that $R_m(x,y)$ iff $R(x,t(y))$, as $R(x,t(y))$ is a 
relation between $x$ and $y$.

\section*{Problem 7}
(a): Consider some recursive function $f(x)$. Hence $f(x)=y$ is recursively enumerable, and thus the relation $R(x,y)$:
\[ \exists z(f(y)=z\wedge x\in\omega_z) \] is recursively enumerable. Thus by the weak recursion theorem, there is some number $n$ such that
$x\in\omega_n$ iff $R(x,n)$ iff $x\in\omega_{f(n)}$.


(b):
1: We can just take $f(x)$ to be $x+1$ (which is recursive, as the relation $x+1=y$ is easily shown to be recursively enumerable), meaning, 
by part (a), there is some number $n$ such that $\omega_n=\omega_{n+1}$.\newline
3: This is a special case of 4, when $f(x)=x$.\newline
4: For a recursive function $f(x)$, $x=f(y)$ is recursively enumerable (because $f(y)=x$ is recursively enumerable). Hence by Problem 5, there is some number $n$ 
such that $x\in\omega_n$ iff $x=f(n)$, that is, $\omega_n=\{f(n)\}$.\newline
5: Consider the relation $y\in\omega_x$, which is recursively enumerable. Hence, by Problem 5, there is some number $n$ such that 
$x\in\omega_n$ iff $n\in\omega_x$.

\section*{Problem 8}
(a): We know that $A\ne\N$ and $A\ne\emptyset$ and that $A$ is a recursive set. Thus there is some number $a\in A$ and some number $b\notin A$. Let $f(x)$ be the function:
\[ f(x)=\begin{cases}
	x\in A: & b \\
	x\notin A: & a \\
	\end{cases}
\]
$f(x)=y$ is equivalent to the relation $R(x,y)$: \[ (x\in A\wedge y=b)\vee (x\notin A\wedge y=a) \] which is recursive as $A$ is recursive. Thus $f(x)$ is recursive. Obviously, 
$x\in A$ iff $f(x)\notin A$.


(b): Consider the $f(x)$ of part (a). Because it is recursive, by Problem 7, there is some number $n$ such that $\omega_{f(n)}=\omega_n$. Because $f(n)\in A$ iff $n\notin A$, 
either 
$f(n)$ or $n$ are in $A$, but not both. But this means that $A$ is not extensional, for otherwise both $f(n)$ and $n$ would need to be in $A$.

\section*{Problem 9}
(a): Take $\phi(i)$ to be the function $d(t(i))$. Thus $x\in\omega_{\phi(i)}$ iff $x\in\omega_{d(t(i))}$ iff $R_{t(i)}(x,t(i))$ iff $M(x,i,\phi(i))$, which is what 
we were required to prove.


\textit{Note to self}: Here is why $M(x,z,d(y))$ is recursively enumerable. Because $M(x,y,z)$ is, then so is $M(x,y,d(z))$. Then let $M_1(x,y,z)$ be the relation 
$M(x,z,d(y))$ which is recursively enumerable. By the iteration theorem, there is some funciton $t(i)$ such that $R_{t(i)}(x,y)$ iff 
$M_1(x,y,i)$ iff $M(x,i,d(y))$.


(b): Take $\phi(i)$ to be the function $t(i,h)$. Thus $x\in\omega_{t(i,h)}$ iff $R_h(x,i,h)$ iff $M(x,i,t(i,h))$, which is what we were required to prove.


\textit{Note to self}: Here is why $R_z(x,y,z)$ can be said to be a recursively enumerable relation. $R_x(y,z,z_1)$ is defined to be a recursively enumerable relation. 
Hence let $R_1(x,y,z)$ be the relation $R_x(y,z,z_1)\wedge x=z_1$, which is recursively enumerable, and is the same as saying $R_z(x,y,z)$.

\section*{Problem 10}
Let $M(x,y,z)$ be $R_y(x,z)$. Then by the strong recursion theorem there exists some function $\phi(i)$ such that 
\[ \omega_{\phi(i)}=\{x:M(x,i,\phi(i))\}=\{x:R_i(x,\phi(i))\} \] which proves Theorem 8 is a special case of the strong recursion theorem.

\section*{Problem 11}
Suppose that $f(x)$ reduces $A$ to $B$ and $\phi(x)$ is a generative function for $A$, and $t(y)$ is an iterative function for the relation $f(x)\in\omega_y$. Hence 
$x\in A$ iff $f(x)\in B$, $x\in\omega_{t(y)}$ iff $f(x)\in\omega_{y}$, and $\phi(i)\in A$ iff $\phi(i)\in\omega_i$. Thus $f(\phi(t(x)))\in B$ iff $\phi(t(x))\in A$ iff $\phi(t(x))\in\omega_{t(x)}$
iff $f(\phi(t(x)))\in\omega_x$. That is, $f(\phi(t(x)))\in B$ iff $f(\phi(t(x)))\in\omega_x$. If we let $\phi(i)$ be the function $f(\phi(t(x)))$, this proves that 
$\phi(i)\in B$ iff $\phi(i)\in\omega_i$.

\section*{Problem 12}
For every $i$, $x\in\omega_i$ iff $\delta(x,i)\in K$. Hence let $f(x)$ be the function $\delta(x,i)$. Thus $x\in\omega_i$ iff $f(x)\in K$, thus showing $K$ is universal under the function 
$f(x)$.


$Note$: Smullyan must have made a typo when defining $K$. He clearly means that $x\in\omega_y$ iff $\delta(x,y)\in K$, but the definition was $\partial(x,y)\in K$, which is an undefined 
function.

\section*{Problem 13}
The set $C$ is recursively enumerable, as it is equivalent to the recursively enumerable relation $x\in\omega_x$. Now consider some universal set $U$. Hence $C$ is reducible to $U$. Because $I(x)$ is a generative function 
for $C$, $f(I(t(x)))$ (as defined in Theorem 9$^*$) is a generative function for $U$ by Theorem $9^*$. 

\section*{Problem 14}
We can show that Post's set $K$ is uniformly universal. We have shown that it is universal under the function $\delta(x,i)$ for every $\omega_i$. Hence $K$ is uniformly universal under the 
recursive function $\delta(x,y)$, as $x\in\omega_i$ iff $\delta(x,i)\in K$. Now consider any universal set $U$. Hence, for some recursive function $f(x)$, 
$x\in K$ iff $f(x)\in U$. Take $f(x,y)$ to be the function $f(\delta(x,y))$. Hence, for all $i$, $f(\delta(x,i))\in U$ iff $\delta(x,i)\in K$ iff $x\in\omega_i$, thus proving that 
$U$ is uniformly universal under the recursive function $f(\delta(x,y))$.

\section*{Exercise 4}
Suppose that $A$ is uniformly universal under $f(x,y)$ and that $t(y)$ is an iterative function for the relation $f(x,x)\in\omega_y$. Thus, for all $i$, $x\in\omega_i$ iff $f(x,i)\in A$, 
and $x\in\omega_{t(y)}$ iff $f(x,x)\in\omega_y$. Hence $f(t(x),t(x))\in A$ iff $t(x)\in\omega_{t(x)}$ iff $f(t(x),t(x))\in\omega_x$, thus proving $f(t(x),t(x))$ is a generative function for 
$A$.

\section*{Problem 15}
(1): We assume that $B$ is a generative set under $\phi(x)$, meaning $\phi(x)\in B$ iff $\phi(x)\in\omega_x$. Thus if $\omega_i=\N$ then, by definition, $\phi(i)\in\omega_i$ (as all 
positive integers are in $\omega_i$) and thus $\phi(i)\in B$.


(2): If $\omega_i=\emptyset$, then $\phi(i)\notin B$. For if it was, then $\phi(i)\in\omega_i$. But this would be contrary to the assumption that $\omega_i=\emptyset$.

\section*{Problem 16}
Suppose that $A$ is a recursively enumerable set. Define $M(x,y)$ iff $y\in A$. Because $M(x,y)$ is recursively enumerable, by the iteration theorem, there is some recursive function 
$t(y)$ such that for all $i$, $\omega_{t(i)}=\{x:M(x,i)\}$. 


(1): Suppose $i\in A$. Hence $x\in\omega_{t(i)}$ iff $M(x,i)$ iff $i\in A$, which shows that, because $i\in A$, $x\in\omega_{t(i)}$ for every $x$.


(2): Suppose $i\notin A$. By part (1) of this problem, we know that $x\in\omega_{t(i)}$ iff $i\in A$. Hence, because $i\notin A$, $x\notin\omega_{t(i)}$ for every $x$, thus meaning 
$\omega_{t(i)}=\emptyset$.

\section*{Problem 17}
Consider some generative set $B$. By Problem 15, we know that, if $B$ is generative under the recursive function $\phi(x)$, then, for every $i$:\newline
(1) If $\omega_i=\N$ then $\phi(i)\in B$.\newline
(2) If $\omega_i=\emptyset$ then $\phi(i)\notin B$.\newline
Now consider any recursively enumerable set $A$. Hence, by Problem 16, there is a recursive function $t(y)$ such that for every number $i$:\newline
(1) If $i\in A$, then $\omega_{t(i)}=\N$;\newline
(2) If $i\notin A$, then $\omega_{t(i)}=\emptyset$.\newline
Now, consider the function $\phi(t(i))$. If $i\in A$, then $\omega_{t(i)}=\N$ and thus $\phi(t(i))\in B$. On the other hand, if $\phi(t(i))\in B$, then 
$\phi(t(i))\in\omega_{t(i)}$. If $i\notin A$, then $\omega_{t(i)}=\emptyset$, which would be contrary to the fact that $\phi(t(i))\in B$ (by (2) of Problem 15). Thus $i\in A$. We have therefore 
shown that $i\in A$ iff $\phi(t(i))\in B$, which shows that $A$ is reducible to $B$ under the function $\phi(t(i))$, and therefore $B$ is universal.
 

\end{document}
